% Copy-ready supplementary extension for the canonical-atom manuscript.
% This file is not included automatically.

\subsection{Enumerating canonical atom attachments}
\label{subsection-canonical-attachment-enumeration}

Fix pairwise disjoint labelled canonical atoms
$A_1,\ldots,A_k$, with every atom vertex labelled, and put
\[
w_i=|V(A_i)|,\qquad
P=\prod_{i=1}^kw_i,\qquad
R=\sum_{i=1}^k(w_i-1).
\]
We count quotient assemblies, not construction histories.  Shared points are
the actual non-singleton vertex classes and carry no external labels.
The atom labels and internal vertex labels are retained.

An admissible connected attachment is an equivalence relation on the disjoint
union of the atom vertex sets such that every non-singleton class contains at
most one vertex from each atom and the atom--shared-point incidence graph is a
tree.  Let $N_t(\mathbf A)$ denote the number of such attachments with exactly
$t$ shared points, and write $(x)_r=x(x-1)\cdots(x-r+1)$.

\begin{theorem}[Exact canonical attachment count]
\label{theorem-exact-canonical-attachment-count}
For $k\ge2$ and $1\le t\le k-1$,
\[
\boxed{
N_t(\mathbf A)
=
\left(\prod_{i=1}^kw_i\right)
S(k-1,t)
\left(\sum_{i=1}^k(w_i-1)\right)_{t-1},
}
\]
where $S(n,t)$ is a Stirling number of the second kind.  The count is zero
outside this range.  For $k=1$, the unique attachment has no shared point.
\end{theorem}

\begin{proof}
Temporarily label the $t$ shared-point nodes.  Write
$x_i=d_i-1$ for the atom excess degrees and
$y_j=\mu_j-1$ for the shared-point excess degrees.  The incidence tree gives
\[
\sum_i x_i=t-1,\qquad
\sum_j y_j=k-1,\qquad
x_i\ge0,\quad y_j\ge1.
\]
The fixed-degree bipartite Prüfer formula, including the injections of incident
point nodes into distinct atom vertices, is
\[
\frac{(t-1)!(k-1)!}
{\prod_i x_i!\prod_j y_j!}
\prod_i(w_i)_{x_i+1}.
\]
Vandermonde's identity gives
\[
\sum_{\sum x_i=t-1}
\prod_i\frac{(w_i)_{x_i+1}}{x_i!}
=
\left(\prod_iw_i\right)\binom{R}{t-1},
\]
while the surjection identity gives
\[
(k-1)!\sum_{\substack{y_j\ge1\\\sum y_j=k-1}}
\frac1{\prod_jy_j!}
=t!S(k-1,t).
\]
The resulting count with labelled point nodes is $t!$ times the displayed
answer.  Forgetting those labels is exactly $t!$-to-one.
\end{proof}

If the quotient has $n$ vertices, then every connected attachment makes
$k-1$ identifications, so $R=n-1$.  Hence
\[
N_t(\mathbf A)
=
\left(\prod_iw_i\right)S(k-1,t)(n-1)_{t-1}.
\tag{10.20}
\]

\begin{corollary}[Profile-refined count]
\label{corollary-profile-refined-attachment-count}
Let $\lambda=(\lambda_1,\ldots,\lambda_t)\vdash k-1$ be the unordered
shared-point excess profile, and let $m_r$ be the multiplicity of the part $r$.
Then the number of attachments with profile $\lambda$ is
\[
\boxed{
\left(\prod_iw_i\right)(R)_{t-1}
\frac{(k-1)!}{\prod_j\lambda_j!\prod_rm_r!}.
}
\]
\end{corollary}

Let $S_{\le q}(n,t)$ be the number of partitions of an $n$-set into $t$
nonempty blocks of size at most $q$.  Summing the preceding corollary gives
\[
N_{t,\le M}(\mathbf A)
=
\left(\prod_iw_i\right)(R)_{t-1}
S_{\le M-1}(k-1,t),
\tag{10.21}
\]
for the number of attachments whose largest shared-point multiplicity is at
most $M$.

\begin{corollary}[Injective-port Cayley identity]
\label{corollary-injective-port-cayley}
For labelled trees on the atom labels,
\[
\boxed{
\sum_{T\in\mathcal T_k}\prod_{i=1}^k(w_i)_{\deg_T(i)}
=
\left(\prod_iw_i\right)(R)_{k-2}.
}
\]
\end{corollary}

Indeed, this is the endpoint $t=k-1$, where all shared points have
multiplicity two and suppress to the edges of an ordinary tree.

The attachment polynomial
\[
\Phi_{\mathbf A}(z)
=
\left(\prod_iw_i\right)
\sum_{t=1}^{k-1}S(k-1,t)(R)_{t-1}z^t
\]
depends on the size vector only through $k$, $\prod_iw_i$, and
$\sum_i(w_i-1)$.  Its positive coefficients are strictly log-concave, hence
unimodal: the Stirling row is strictly log-concave by the real-rootedness of
the Touchard polynomial, and the falling-factorial row is strictly
log-concave because
\[
\frac{(R)_{t-1}^2}{(R)_{t-2}(R)_t}
=
\frac{R-t+2}{R-t+1}>1.
\]

For fixed $k$ and $R\to\infty$, the binary endpoint dominates.  More precisely,
\[
\frac{N_{k-2}}{N_{k-1}}
=
\frac{\binom{k-1}{2}}{R-k+3},
\]
and all lower terms are $O_k(R^{-2})$ relative to $N_{k-1}$.
